\section{H-Net: convertir la tokenización en una jerarquía aprendible}

El capítulo anterior señaló que el propósito de la tokenización no es solo acortar las secuencias, sino crear unidades sobre las que puede actuar eficazmente la atención. El objetivo de diseño de H-Net no es simplemente cambiar un tokenizer, sino integrar la segmentación, la compresión y el modelado de secuencias en una misma red end-to-end, permitiendo que el modelo aprenda recursivamente múltiples niveles de abstracción. Además, explica cómo Transformer y SSM pueden asumir responsabilidades distintas a diferentes resoluciones dentro de la misma arquitectura.

\subsection{Dynamic chunking: los límites son predicciones dependientes del contexto}

H-Net comienza desde unidades de bajo nivel como el byte; un encoder genera sus representaciones y luego predice en qué posiciones se ubican los límites de los chunks. Los límites verdes/coloridos a la derecha de la figura exploran activamente durante las primeras etapas del entrenamiento, alineándose gradualmente más tarde con espacios, subwords y combinaciones de niveles superiores; no son etiquetas manuales ni garantizan ser equivalentes a tokens lingüísticos, sino que se moldean conjuntamente por el loss de modelado de lenguaje en decisiones de compresión.

\lecturefigure{slide-025.jpg}{H-Net: aprender chunks dinámicos dependientes de los datos desde una secuencia raw}{Grabación oficial de Stanford, 00:34:10--00:37:19}

Sea $z_1,\ldots,z_L$ la salida del encoder y sea el head de enrutamiento responsable de las probabilidades de los límites:
\[
r_t=\sigma(w^\top z_t+b),
\qquad
m_t\in\{0,1\},\quad m_t\sim\operatorname{Route}(r_t).
\]
Donde $r_t$ representa la tendencia de la posición $t$ para ser el final del chunk y $m_t$ es la decisión de enrutamiento. Si el chunk $j$ cubre el intervalo $I_j$, su resumen puede expresarse como
\[
c_j=\operatorname{Pool}\{z_t:t\in I_j\}.
\]
El H-Net real requiere un enrutamiento diferenciable, suavizado, normalización y optimización consciente de las etapas; la ecuación anterior solo muestra la interfaz para "predecir límites y agregar intervalos".

\teachervoice{00:35:07--00:37:17, el orador describe una animación del entrenamiento: inicialmente el modelo ignora los límites, luego se estabilizan cerca de espacios y subwords con significado; las combinaciones multinivel pueden formar grupos parecidos a frases. Este fenómeno es evidencia de compresión aprendida, no de la verdad de fondo (ground truth) del tokenizer.}

\subsection{Arquitectura completa: chunking, modelo principal y deschunking}

La figura izquierda muestra una jerarquía en forma de ampolla: una secuencia de grano fino pasa primero por un encoder externo y el proceso de chunking para acortarse; el modelo principal opera luego sobre unidades más gruesas y finalmente el deschunking expande las representaciones a la resolución original antes de que el decoder realice predicciones autoregresivas. La figura derecha despliega el módulo de enrutamiento y suavizado. Al observar esta imagen, cabe destacar que el modelo principal puede ser un Transformer; H-Net no implica "deshabilitar la atención en toda la red".

\lecturefigure{slide-026.jpg}{Arquitectura de H-Net: chunking multinivel, modelo principal, deschunking y estabilización de señales}{Grabación oficial de Stanford, 00:37:20--00:40:39}

\begin{lstlisting}[caption={Flujo de datos didáctico para una sola jerarquía de H-Net}]
def hnet_stage(fine_sequence, outer_encoder, router, main_model, decoder):
    fine_features = outer_encoder(fine_sequence)
    boundary_scores = router(fine_features)
    coarse_features, routing_map = chunk(fine_features, boundary_scores)
    coarse_outputs = main_model(coarse_features)
    expanded_outputs = dechunk(coarse_outputs, routing_map)
    return decoder(fine_features, expanded_outputs)
\end{lstlisting}

\begin{knowledgebox}{¿Por qué la etapa externa tiende a SSM mientras que la interna puede usar Transformer}
El encoder externo se enfrenta directamente a unidades de alta resolución y débil semántica como bytes o caracteres, necesitando una compresión en línea para formar chunks; un estado recurrente limitado posee el sesgo inductivo adecuado para esta tarea. El modelo interno recibe ya chunks más escasos y de nivel superior, por lo que la atención por chunk se vuelve razonable. La arquitectura no es una elección binaria, sino hacer que los primitivos coincidan con la resolución.
\end{knowledgebox}

\begin{warningbox}{El enrutamiento dinámico es un problema de optimización discreta difícil}
La decisión de los límites altera la longitud de la secuencia y las rutas de cálculo posteriores, lo que hace propenso al entrenamiento a cortes totales, sin cortes, inestabilidad de gradientes o colapso de etapas. El orador afirma explícitamente que el H-Net publicado es "el primer paso funcional", dependiendo del módulo de enrutamiento, suavizado, normalización, proyección y optimización consciente de las etapas; no se debe inferir una implementación sencilla a partir de estas cuatro líneas de pseudocódigo.
\end{warningbox}

\teachervoice{00:40:00--00:40:43, Gu menciona que esto es una optimización discreta no trivial; en la Q\&A de 00:57:23--00:59:08 se enfatiza nuevamente que H-Net aún no ha sido validado completamente a gran escala, y tanto el mecanismo de chunking como cada etapa tienen mucho margen de mejora.}

\subsection{Resumen de la sección}

H-Net reescribe un pipeline offline de tokenizer--LM--detokenizer como una jerarquía aprendible: los encoders de bajo nivel forman chunks dependientes del contexto, el modelo interno realiza cálculos sobre unidades de grano grueso y luego se expande a la resolución original para predecir. Su aportación sistémica más importante es la división del trabajo por resolución: SSM asume la compresión en los límites finos originales, mientras que la atención asume la combinación precisa sobre chunks de alto nivel. El próximo capítulo utilizará scaling y ablation para verificar si este diseño realmente genera un mejor sesgo inductivo.


